% from ega3-7-fr.tex
\oldpage[III]{178}
En particulier, si $B$ est une $A$-algèbre (non nécessairement commutative),
on a un homomorphisme fonctoriel en $M$
\[
  (T(M))_{(B)}
  =
  T(M)\otimes_A B
  \longrightarrow
  T(M\otimes_A B)
  =
  T_{(B)}(M_{(B)})
  \tag{7.2.3.2}\label{III.7.2.3.2.fr}
\]
qui, en vertu de la fonctorialité de (7.2.3.1) en $M$, est un homomorphisme
de $B$-modules.

On a en outre le diagramme commutatif
\[
\xymatrix@C=30pt@R=30pt{
  T(A)\otimes_A M
  \ar[r]^-{t_M}\ar[d]
  &
  T(M)\ar[d]
  \\
  T_{(B)}(B_s)\otimes_B M_{(B)}
  \ar[r]_-{t_{M_{(B)}}}
  &
  T_{(B)}(M_{(B)})
}
\tag{7.2.3.3}\label{III.7.2.3.3.fr}
\]
où la flèche verticale de droite est l'homomorphisme composé
\[
  T(M)
  \longrightarrow
  T(M)\otimes_A B
  \longrightarrow
  T(M\otimes_A B)
  =
  T_{(B)}(M_{(B)})
\]
de (7.2.3.2) et de l'homomorphisme canonique; quant à la flèche verticale de
gauche de (7.2.3.3), c'est l'homomorphisme
\[
  T(A)\otimes_A M
  \longrightarrow
  T_{(B)}(B_s)\otimes_B(B\otimes_A M)
  =
  T_{(B)}(B_s)\otimes_A M
\]
où
\[
  T(A)\longrightarrow T_{(B)}(B_s)=T(B)
\]
est $T(\rho)$, $\rho$ étant considéré comme homomorphisme de $A$-modules
$A\to B_{[\rho]}$.
\end{env}

\begin{lemma}[7.2.4]
Si $T$ est un foncteur covariant additif de $\mathbf{Ab}_A$ dans
$\mathbf{Ab}$, commutant avec les sommes directes, l'homomorphisme canonique
$t_L$ (7.2.2.1) est un isomorphisme pour tout $A$-module libre $L$.
\end{lemma}

\begin{proof}
En effet, on a
\[
  L=\bigoplus_{\alpha\in I}L_\alpha
\]
où $L_\alpha$ est isomorphe à $A_s$ pour tout $\alpha\in I$; la définition de
$t_M$ donnée dans (7.2.2) montre que
\[
  t_L=\bigoplus_{\alpha\in I}t_{L_\alpha},
\]
puisque
\[
  T:
  \operatorname{Hom}_A(A_s,L)
  \longrightarrow
  \operatorname{Hom}_{\mathbf Z}(T(A_s),T(L))
\]
est somme directe des applications $\mathbf Z$-linéaires
\[
  T_\alpha:
  \operatorname{Hom}_A(A_s,L_\alpha)
  \longrightarrow
  \operatorname{Hom}_{\mathbf Z}(T(A_s),T(L_\alpha))
\]
en vertu de l'hypothèse sur $T$. On est donc ramené à prouver le lemme pour
$L=A_s$; mais $t_L$ n'est autre alors que l'isomorphisme canonique
\[
  T(A_s)\otimes_A A_s\xrightarrow{\sim}T(A_s)
\]
valable pour tout $A$-module à droite.
\end{proof}

\begin{proposition}[7.2.5]
Soit $T$ un foncteur covariant additif de $\mathbf{Ab}_A$ dans
$\mathbf{Ab}$, commutant aux sommes directes. Les conditions suivantes sont
équivalentes :
\begin{enumerate}
  \item[a)] $T$ est exact à droite.
  \item[b)] L'homomorphisme canonique $t_M$ (7.2.2.1) est un isomorphisme pour
    tout $A$-module à gauche $M$.
  \item[$b'$)] $T$ est semi-exact et l'homomorphisme $t_M$ est surjectif pour
    tout $A$-module à gauche $M$.
  \item[c)] $T$ est isomorphe à un foncteur en $M$ de la forme
    $N\otimes_A M$, où $N$ est un $A$-module à droite.
\end{enumerate}
\end{proposition}

\begin{proof}
Il est clair que $b)$ implique $c)$ et que $c)$ implique $a)$; montrons que
$a)$ implique $b)$.
